\begin{recetaConImagen}{Mousse de Limón}{5-6 vasitos}{10 min activos · 2 h total}{receta:mousse_limon}
    {imagenes/mousse_limon.png}{Fresquita, cremosa y con ese toque de limón que entra solo.}

    \begin{minipage}[t]{0.35\textwidth}
        \section*{Ingredientes}
        \begin{itemize}[leftmargin=*]
            \item 500 g de yogur griego
            \item 1 bote de leche condensada (aprox. 370 g)
            \item el zumo de 1 limón grande
            \item ralladura de limón
        \end{itemize}
    \end{minipage}
    \hfill
    \begin{minipage}[t]{0.6\textwidth}
        \section*{Preparación}
        \begin{enumerate}
            \item \textbf{Mezcla base:} Poner en la batidora el yogur griego, la leche condensada y el zumo del limón.
            \item \textbf{Batir:} Triturar todo hasta que coja una textura suave, homogénea y con cuerpo.
            \item \textbf{Servir:} Repartir la mousse en recipientes individuales.
            \item \textbf{Toque final:} Añadir por encima un poco de ralladura de cáscara de limón.
            \item \textbf{Frío:} Dejar enfriar en la nevera hasta el momento de servir.
        \end{enumerate}
    \end{minipage}

    \vspace{1em}
    \begin{tcolorbox}[colback=orange!5, colframe=orange!50, title=Consejo, size=small]
        Si quieres una mousse más fresca e intensa, añade la ralladura de limón justo antes de servir para que conserve mejor su aroma.
    \end{tcolorbox}
\end{recetaConImagen}
